\begin{helper}
Let \(V\) be finite, let \(G\) be a graph on \(V\), and suppose
\(\gamma>0\) and \(\gamma\le\Delta\).  For the Bernoulli activation and
uniform-priority law constructed in
\(\noderef{SamplingFiniteProductPriorityLawExists}\), write
\[
  p=\gamma/\Delta .
\]
Then the same sampling law has the following finite-coordinate cylinder
extension property.  For any two disjoint finite coordinate sets \(C,D\),
activation atoms \(A_0\subseteq C\) and \(A_D\subseteq D\), lower-cut sets
\(L\subseteq C\), \(K\subseteq D\), upper-cut set \(U\subseteq C\), and
thresholds \(t_v\in[0,1]\) on the used coordinates, provided the lower and
upper cut sets are disjoint in the indicated sense, there are nonnegative
real masses \(w_{\rm base}\) and \(w_{\rm ext}\) such that
\[
\operatorname{ofReal}(w_{\rm base})
 = \nu\{(A,\pi):A\cap C=A_0,\ \pi_v\le t_v\ (v\in L),\
        t_v<\pi_v\ (v\in U)\},
\]
\[
\operatorname{ofReal}(w_{\rm ext})
 = \nu\{(A,\pi):A\cap(C\cup D)=A_0\cup A_D,\
        \pi_v\le t_v\ (v\in L\cup K),\
        t_v<\pi_v\ (v\in U)\},
\]
and
\[
  w_{\rm ext}
  =
  w_{\rm base}\,
  p^{|A_D|}(1-p)^{|D|-|A_D|}
  \prod_{v\in V}\begin{cases}t_v,&v\in K,\\1,&v\notin K.\end{cases}
\]
The law also retains its total mass, activation-atom formula, and full
priority lower-rectangle formula.
\end{helper}
\begin{proof}
Apply \(\noderef{SamplingFiniteProductPriorityLawExists}\) to
\(G,\Delta,\gamma\).  Its statement gives a probability measure \(\nu\) on
activation-priority samples with total mass \(1\), the Bernoulli activation
atom formula with parameter \(p=\gamma/\Delta\), and the full priority
lower-rectangle formula.

The inequalities \(\gamma>0\) and \(\gamma\le\Delta\) imply
\(0\le p\le 1\): nonnegativity follows because both numerator and denominator
are nonnegative, and the upper bound follows by dividing the inequality
\(\gamma\le\Delta\) by the positive real \(\Delta\).  Therefore all
hypotheses of \(\noderef{SamplingFiniteCoordinateCellExtensionMass}\) are
satisfied for the chosen measure \(\nu\), the parameter \(p\), the disjoint
coordinate sets \(C,D\), the activation atoms \(A_0,A_D\), the lower and
upper cut sets \(L,K,U\), and the threshold function \(t\).  That node states
exactly the existence of nonnegative representatives \(w_{\rm base}\) and
\(w_{\rm ext}\) for the displayed base and extended finite-coordinate cells,
with
\[
  w_{\rm ext}
  =
  w_{\rm base}\,
  p^{|A_D|}(1-p)^{|D|-|A_D|}
  \prod_{v\in V}\begin{cases}t_v,&v\in K,\\1,&v\notin K.\end{cases}.
\]
Substituting \(p=\gamma/\Delta\) gives the asserted formula, while the total
mass and the two law clauses are the clauses already supplied by
\(\noderef{SamplingFiniteProductPriorityLawExists}\).
\end{proof}
